
@@M0@@ qui est représentable par un schéma affine sur @@M1@@. Or @@M2@@ est défini par un faisceau inversible @@M3@@ sur @@M4@@ muni d'une rigidification @@M5@@ pour @@M6@@. De même, un point @@M7@@ correspond à un faisceau inversible @@M8@@ sur @@M9@@ muni d'une rigidification @@M10@@. L'image de @@M11@@ dans @@M12@@ est @@M13@@ si et seulement s'il existe un isomorphisme de faisceaux rigidifiés @@M14@@ où @@M15@@ est l'inclusion pour @@M16@@. Remarquons que ces isomorphismes définissent un isomorphisme @@M17@@ tel que @@M18@@.

Réciproquement si @@M19@@ est un isomorphisme @@M20@@ tel que @@M21@@, il existe un faisceau inversible rigidifié @@M22@@ sur @@M23@@ qui s'obtient par recollement suivant @@M24@@~; explicitement @@M25@@ est le noyau de l'homomorphisme canonique @@M26@@ défini par @@M27@@. Ce @@M28@@ correspond alors à un point @@M29@@ de @@M30@@ qui va sur @@M31@@.

Finalement soit @@M32@@ le foncteur défini, lorsque @@M33@@ parcourt les schémas sur @@M34@@, par les ensembles d'isomorphismes @@M35@@ tels que @@M36@@. On vient d'établir un isomorphisme entre l'image réciproque de @@M37@@ dans @@M38@@ et le foncteur @@M39@@. Il suffit donc d'appliquer le lemme suivant (où l'on prend @@M40@@)~:

\textbf{Lemme 2.4.} \uline{Soit} @@M41@@ \uline{un schéma plat, propre et de présentation finie sur} @@M42@@\uline{. Soient} @@M43@@ \uline{deux faisceaux inversibles sur} @@M44@@, @@M45@@ \uline{le foncteur des isomorphismes} @@M46@@ \uline{(pour} @@M47@@\uline{). Alors}
\begin{enumerate}[label=\roman*)]
\item @@M48@@ \uline{est représentable par un schéma affine sur} @@M49@@ \uline{de présentation finie.}
\item \uline{Si de plus} @@M50@@ \uline{et} @@M51@@ \uline{sont munis de rigidifications (relativement à une}
